(1) 取正交横向基 $\hat x,\hat y$，右、左圆偏振基向量可记为 $\mathbf e_R=(\hat x-i\hat y)/\sqrt2$、$\mathbf e_L=(\hat x+i\hat y)/\sqrt2$。入射沿 $\hat x$ 的线偏振电场分解为 $\hat x=(\mathbf e_R+\mathbf e_L)/\sqrt2$，两圆分量振幅相同。
(2) 两圆偏振穿过厚度 $L$ 分别积累相位 $\phi_R=2\pi n_RL/\lambda_0$、$\phi_L=2\pi n_LL/\lambda_0$。重新组合后整体公共相位不改变偏振方向，而相对相位 $\Delta\phi=\phi_R-\phi_L$ 使线偏振方位旋转 $\alpha=\Delta\phi/2=\pi(n_R-n_L)L/\lambda_0$，符号按圆偏振约定确定。
(3) 因此旋转角与厚度成正比，在 $n_R-n_L$ 不随波长变化时按 $1/\lambda_0$ 改变；真实旋光材料还可具有明显色散，必须写成 $\alpha(\lambda_0)=\pi L[n_R(\lambda_0)-n_L(\lambda_0)]/\lambda_0$。若 $n_R=n_L$，相对相位为零，输出偏振方向不旋转。
(4) 若左右圆偏振具有不同吸收系数，则两分量在出射处的幅值不再相等，除了相对相位还出现幅度差。两个不等幅圆分量叠加通常形成椭圆偏振，因此有圆二色性时仅报告旋转角不足以完整描述光的偏振态。 因此实验上除了分析偏振椭圆的长轴取向，还应测量其椭圆率或Stokes参数，才能区分圆双折射导致的旋转和圆二色性引起的椭圆化。 具体地，归一化圆偏振Stokes参数 $s_3=S_3/S_0$ 可揭示两个圆偏振分量之间的振幅不平衡。
